\begin{abstract}
\textbf{Objective:} We examine whether the apparent superiority of
discrimination-oriented EEG regularisation survives correcting
the covariance criterion used in held-source validation.
\textbf{Methods:} We audit the original zero-target covariance/
tangent-space linear-discriminant experiments and distinguish
a validation-class-centred Gaussian covariance log score
from class-conditional predictive log scoring under source-
training class means. An exact nonnegative mean-shift
penalty separates the two estimands. Original prospective
participant-held experiments are preserved. Proper-predictive
comparators are subsequently executed on two already
exposed, separately archived public EEG datasets and
are classified as post-outcome analyses.
\textbf{Results:} The original prospective EEGMMIDB
block ($N=20$) showed $+0.0561$ mean AUC for
source-discrimination selection over the historical
class-centred covariance criterion (Holm $p=0.00597$);
five participants were materially harmed. Two original
external BNCI tests were individually non-significant,
and a separate prospective filter-bank hypothesis was
refuted. On retrospective Shin2017A ($N=29$),
source-discrimination, proper predictive and historical
class-centred selection obtained AUCs of 0.785900,
0.786207 and 0.763678, respectively.
On retrospective BNCI2014\_002 ($N=14$), the
corresponding AUCs were 0.778705, 0.771886 and
0.761116. The proper-minus-class-centred
source Gaussian log-score penalty averaged 0.86995
and 0.91218 per feature in the two datasets,
remaining positive across all 43 overlapping
outer training pools.
\textbf{Conclusion:} Historical class-centred covariance
validation is not proper source-mean predictive
likelihood. Correcting the estimand substantially
narrows the claimed discrimination-selector
advantage without establishing a universally
superior shrinkage rule. The retrospective comparisons
require genuinely untouched external confirmation.
\end{abstract}
